\subsection{The resonant sign from the original upper field}
\label{rw:app:upper-resonant-sign}

The incoming lower-loss chart uses the sign of the first orbital
resonant coefficient of the original upper field. We compute it from
\eqref{rw:upper:divided-field}, beginning on $\nu=0$.
Its equilibrium satisfies
\[
 X_s=-\bar cY_s,\qquad
 Y_s=\frac{1+\epsilon\bar a\bar c}
           {\bar m+\epsilon\bar B\bar c^2}.
\]
The trace is zero precisely when
\begin{equation}\label{rw:finite:zero-trace}
 0=\bar c[1+\epsilon(\bar m-2\bar B)]
 +\epsilon\bar a\bar m
 +\epsilon\bar a\bar c^2(1-\epsilon\bar B)
 +\epsilon^2\bar B\bar c^3.
\end{equation}
Indeed the trace, multiplied by $Y_s+\epsilon$, is
$\bar cY_s(1-2\epsilon\bar B)+\epsilon(\bar a+\bar c)$;
substitution of $Y_s$ gives the displayed polynomial. The implicit
function theorem yields
$\bar c=-\epsilon\bar a\bar m+O(\epsilon^2)$.
Thus $X_s=\epsilon\bar a+O(\epsilon^2)$ and
$Y_s=1/\bar m+O(\epsilon^2)$.

Put $x=X-X_s$, $y=Y-Y_s$. To first order in $\epsilon$,
expansion of $1/(Y+\epsilon)$ in the first component gives the
quadratic and cubic terms
\begin{align*}
 \epsilon P_2&=\epsilon(\bar B\bar mx^2
               -\bar a\bar m^2xy+\bar m^3y^2),\\
 \epsilon P_3&=\epsilon(-\bar B\bar m^2x^2y
               +\bar a\bar m^3xy^2-\bar m^4y^3).
\end{align*}
The second component is affine, and the omitted coefficients are
$O(\epsilon^2)$ in the reserved finite coefficient jets.
At zero scale use $x=\sqrt{\bar m}(u-v)$, $y=u+v$.
The linear part is $(\sqrt{\bar m}u,-\sqrt{\bar m}v)$.
The coefficients of $u^2v$ and $uv^2$ in $P_3$ are respectively
\[
 \bar B\bar m^3+\bar a\bar m^{7/2}-3\bar m^4,
 \qquad
 \bar B\bar m^3-\bar a\bar m^{7/2}-3\bar m^4.
\]
Since $\dot u=(\dot y+\dot x/\sqrt{\bar m})/2$ and
$\dot v=(\dot y-\dot x/\sqrt{\bar m})/2$, the sum of the two
resonant vector-field coefficients is
$\epsilon\bar a\bar m^3+O(\epsilon^2)$.
The quadratic terms affect this coefficient only at
$O(\epsilon^2)$: their coefficients and normalizing changes are
$O(\epsilon)$, since the nonresonant denominators have fixed margins.
The $O(\epsilon)$ change of linear eigenvectors applied to $P_3$
has the same order of effect.

Dividing the two components to form the orbital equation changes
this sum by the factor $-1/\lambda_u$. Thus its $V^2$ coefficient
in the notation of \eqref{rw:upper:axis-equation} is
$-\epsilon\bar a\bar m^{5/2}+O(\epsilon^2)$.
The analytic construction of Section~\ref{rw:upper-resonance}
vanishes exactly both at $\epsilon=0$ and on the reversible face.
Restricting to $\kappa=\nu=0$ and applying these two divisions
therefore strengthens the remainder to
\begin{equation}\label{rw:finite:resonant-sign-formula}
 \sigma=-\chi_{\rm loc}^2\epsilon\bar a
                 [\bar m^{5/2}+O(\epsilon)].
\end{equation}
Here $\chi_{\rm loc}>0$ is the common local section scale.
The two axis shears have no quadratic terms affecting this constant
coefficient; the common section rescaling multiplies it by
$\chi_{\rm loc}^2$. The sign is unchanged by the common reversal
$U=-u,V=-v$ that selects the returning quadrant.
On a compact block $\bar a>a_*>0$, analyticity of $\sigma/\epsilon$
then gives $\sigma<-c\epsilon$ for small
$|\kappa|+|\nu|+\epsilon$. This supplies the sign for the prepared
original field on that compact block.
